\newpage
\section*{РЕЗЮМЕ СТАРТАП-ПРОЕКТА}
\addcontentsline{toc}{section}{РЕЗЮМЕ СТАРТАП-ПРОЕКТА}


Название: «Информационно-аналитическая система автоматизации проверки и оформления учебно-методических материалов редакционно-издательского отдела». Проект представляет собой web-ориентированную систему, предназначенную для автоматизации процессов проверки, оформления и обработки учебно-методических документов в редакционно-издательском отделе образовательной организации.

Система обеспечивает автоматическую проверку DOCX-документов на соответствие требованиям оформления, генерацию титульных листов, формирование PDF-отчётов о найденных ошибках, а также централизованное хранение и обработку результатов проверки.

Ключевые возможности системы включают:
\begin{enumerate}
	\item Автоматическую проверку учебно-методических материалов на соответствие установленным требованиям оформления с выделением найденных ошибок.
	
	\item Генерацию титульных листов для методических указаний, учебных пособий и монографий на основе пользовательских данных.
	
	\item Формирование PDF-отчётов с результатами проверки документов и статистикой найденных нарушений.
	
	\item Реализацию системы аутентификации и разграничения прав доступа пользователей.
	
	\item Ведение статистики проверок и централизованное хранение информации о документах и пользователях.
	
	\item Интуитивно понятный web-интерфейс, обеспечивающий удобную работу пользователей с системой.
\end{enumerate}

Разработанная система позволяет существенно сократить время проверки учебно-методических материалов, уменьшить количество ошибок оформления и повысить эффективность работы редакционно-издательского отдела.

Актуальность проекта обусловлена необходимостью автоматизации процессов проверки и оформления документов в образовательных организациях. Традиционная ручная проверка требует значительных временных затрат и сопровождается высокой вероятностью появления ошибок, связанных с человеческим фактором.

Целью разработки является создание информационно-аналитической системы, обеспечивающей:
\begin{itemize}
	\item автоматизированную проверку учебно-методических материалов;
	
	\item автоматическое формирование титульных листов и сопроводительной документации;
	
	\item повышение точности проверки документов;
	
	\item сокращение времени обработки учебных изданий;
	
	\item централизованное управление процессами проверки и хранения документов.
\end{itemize}

Отличительные особенности разработанной системы по сравнению с существующими решениями:
\begin{enumerate}
	\item Комплексная автоматизация процессов проверки и оформления учебно-методических материалов в рамках единой системы.
	
	\item Автоматическое формирование PDF-отчётов с визуализацией найденных ошибок.
	
	\item Поддержка генерации титульных листов для различных типов учебных изданий.
	
	\item Использование клиент-серверной архитектуры, обеспечивающей возможность масштабирования системы.
	
	\item Гибкая система управления пользователями и разграничения прав доступа.
\end{enumerate}

В системе реализованы функции автоматической проверки DOCX-документов, генерации титульных листов, формирования PDF-отчётов, ведения статистики проверок, а также механизмы авторизации и взаимодействия клиентской и серверной части через REST API.

Идея создания проекта возникла в 2025 году в связи с необходимостью автоматизации процессов проверки учебно-методических материалов в редакционно-издательском отделе образовательной организации. Разработка системы была начата в 2026 году в рамках выполнения выпускной квалификационной работы.

Бизнес-цели:
\begin{enumerate}
	\item Снижение временных затрат на проверку и оформление учебно-методических материалов.
	
	\item Повышение качества и единообразия оформления документов.
	
	\item Сокращение количества ошибок, связанных с ручной проверкой документов.
	
	\item Повышение эффективности взаимодействия сотрудников редакционно-издательского отдела и преподавателей.
	
	\item Возможность дальнейшего расширения функциональности системы и интеграции с внутренними информационными системами образовательной организации.
\end{enumerate}